% ============ 第2章 §2.13 有界线性泛函(原书 pp.155-158) ============
\section{有界线性泛函}\label{sec:2.13}
\index{圆板代数}


当 $0 < p \leq \infty$ 时,对固定的 $a \in D$,将 $f(a)$ 对应于 $f \in H^{p}$ 的映射,由引理~\ref{lem:2.10.1} 是有界线性泛函.本节考察 $H^{p}$ 上有界线性泛函的全体 $(H^{p})^{*}$.

由于 $H^{p}$ 与其边界值的全体 $H_{*}^{p}$ 由定理~\ref{thm:2.7.1} 与定理~\ref{thm:2.10.5} 可视为同一,利用 $H^{p}$ 是 $L^{p}$ 的线性子空间这一事实来考察 $(H^{p})^{*}$.当 $1 \leq p < \infty$ 时,设 $1/p + 1/q = 1$,则 $(L^{p})^{*}$ 与 $L^{q}$ 等距同构,即 $(L^{p})^{*} = L^{q}$ 是众所周知的.另外当 $0 < p < 1$ 时,$(L^{p})^{*} = \{0\}$ 也是已知的.但如上所述,一般 $(H^{p})^{*} \neq \{0\}$.

设 $1 \leq p \leq \infty$,则因 $H^{p} \subset L^{p}$,把 $L^{p}$ 上的有界线性泛函限制到 $H^{p}$ 上即得 $(H^{p})^{*}$ 的元素.又由 Hahn--Banach 定理,$H^{p}$ 上的有界线性泛函可扩张到 $L^{p}$ 上,故得 $(L^{p})^{*} = L^{q}$($p \neq \infty$)的元素.将这一事实准确叙述即为定理~\ref{thm:2.13.1}.

\begin{theorem}\label{thm:2.13.1}
设 $1 \leq p < \infty$ 且 $1/p + 1/q = 1$.$(H^{p})^{*}$ 与 $L^{q}/zH^{q}$ 等距同构.即 $(H^{p})^{*} = L^{q}/zH^{q}$.
\end{theorem}

\begin{proof}
对 $g \in L^{q}$,令
\begin{equation*}
\Phi_{g}(k) = \int_{-\pi}^{\pi} k(e^{i\theta}) g(e^{i\theta})\dd \theta / 2\pi \quad (k \in L^{p}),
\end{equation*}
则 $\Phi_{g} \in (L^{p})^{*}$ 且 $\|\Phi_{g}\| = \|g\|_{q}$.令 $\phi = \Phi_{g}|H^{p}$,则 $\phi \in (H^{p})^{*}$.由定理~\ref{thm:2.10.5},
\begin{align*}
(H^{p})^{\perp} & = \left\{g \in L^{q}; \int_{-\pi}^{\pi} f(e^{i\theta}) g(e^{i\theta})\dd \theta / 2\pi = 0 \ (f \in H^{p}) \right\} \\
& = zH^{q},
\end{align*}
因此对任意 $h \in zH^{q}$,$\phi = \Phi_{g+h}|H^{p}$.于是 $L^{q}/zH^{q}$ 可缩小地嵌入 $(H^{p})^{*}$.等距性由下面的逆命题的证明得出.

反过来,设 $\phi \in (H^{p})^{*}$.由 Hahn--Banach 定理,存在 $\Phi \in (L^{p})^{*}$,使得 $\phi = \Phi|H^{p}$ 且 $\|\phi\| = \|\Phi\|$.由 $(L^{p})^{*} = L^{q}$,存在 $g \in L^{q}$,使得 $\Phi = \Phi_{g}$ 且 $\|\Phi\| = \|g\|_{q}$.对任意 $h \in zH^{q}$,$\Phi|H^{p} = \Phi_{g+h}|H^{p}$,故 $\|g\|_{q} = \inf \|g + h\|_{q}$.于是 $(H^{p})^{*}$ 可等距地嵌入 $L^{q}/zH^{q}$.
\end{proof}

当 $p = \infty$ 时,因 $(L^{\infty})^{*} \supsetneq L^{1}$,定理~\ref{thm:2.13.1} 不成立.但类似于定理~\ref{thm:2.13.1} 的证明,可以证明 $(H^{\infty})^{*} = (L^{\infty})^{*} / (H^{\infty})^{\perp}$.

\begin{theorem}\label{thm:2.13.2}
设 $1 < p \leq \infty$ 且 $1/p + 1/q = 1$.$H^{p}$ 与 $(L^{q}/zH^{q})^{*}$ 等距同构.即 $H^{p} = (L^{q}/zH^{q})^{*}$.
\end{theorem}

\begin{proof}
对 $g \in H^{p}$,令
\begin{equation*}
\Phi_{g}(k) = \int_{-\pi}^{\pi} k(e^{i\theta}) g(e^{i\theta})\dd \theta / 2\pi \quad (k \in L^{q}),
\end{equation*}
则 $\Phi_{g} \in (L^{q})^{*}$ 且 $\|\Phi_{g}\| = \|g\|_{p}$.因 $\Phi_{g}$ 在 $zH^{q} = (H^{p})^{\perp}$ 上为零,令 $\phi_{g}(k + zH^{q}) = \Phi_{g}(k)$,则 $\phi_{g} \in (L^{q}/zH^{q})^{*}$ 且 $\|\phi_{g}\| \leq \|g\|_{p}$.于是 $H^{p}$ 可缩小地嵌入 $(L^{q}/zH^{q})^{*}$.

反过来,设 $\phi \in (L^{q}/zH^{q})^{*}$,令
\begin{equation*}
\Phi(k) = \phi(k + zH^{q}) \quad (k \in L^{q}),
\end{equation*}
则 $\Phi \in (L^{q})^{*}$ 且 $\|\Phi\| \leq \|\phi\|$.由 $(L^{q})^{*} = L^{p}$,存在 $g \in L^{p}$,使得 $\Phi = \Phi_{g}$ 且 $\|\Phi\| = \|g\|_{p}$.因 $\Phi_{g}$ 在 $zH^{q}$ 上为零,$g \in (zH^{q})^{\perp} = H^{p}$.于是 $(L^{q}/zH^{q})^{*}$ 可缩小地嵌入 $H^{p}$.
\end{proof}

定理~\ref{thm:2.13.2} 在 p = 1 时不成立,但利用 Riesz 兄弟定理,$H^{1}$ 成为某个 Banach 空间的共轭空间.

令
\begin{equation*}
C = \{f \in C[-\pi, \pi]; f(-\pi) = f(\pi)\},
\end{equation*}
则 $C$ 可与 $D$ 的边界 $\partial D$ 上复值连续函数的全体视为同一.

令
\begin{equation*}
A = \{f \in C; \hat{f}(n) = 0 \ (n < 0)\},
\end{equation*}
则其与 $H(D)$ 中在 $\overline{D}$ 上连续的函数限制到 $\partial D$ 上所得者一致.$A$ 称为圆盘代数.

\begin{theorem}\label{thm:2.13.3}
$H^{1}$ 与 $(C/zA)^{*}$ 等距同构.即 $H^{1} = (C/zA)^{*}$.
\end{theorem}

\begin{proof}
因 $C^{*}$ 与 $\partial D$ 上(有界)复测度的全体 $M(\partial D)$ 一致,由定理~\ref{thm:2.12.2},
\begin{align*}
(zA)^{\perp} & = \{\mu \in M(\partial D); \hat{\mu}(n) = 0 \quad (n < 0)\} \\
& = H^{1},
\end{align*}
因此,定理可以与引理 2.13.2 的证明完全相同地证明.
\end{proof}
